% `tex.web` §851, §854 and §873: what the *final* pass may do that no other
% pass may.
%
%     if (b>inf_bad)or(pi=eject_penalty) then
%       begin if final_pass and (minimum_demerits=awful_bad) and
%          (link(r)=last_active) and (prev_r=active) then
%         artificial_demerits:=true {set demerits zero, this break is forced}
%       else if b>threshold then goto deactivate;
%       node_r_stays_active:=false;
%       end
%
% Two things at once. The break is *forced* — taken however bad it is, rather
% than the paragraph failing — and it is charged *nothing*, which §856 prints
% as `d=*`. §852 gives an overfull line `b := inf_bad+1`, and §856 prints
% anything past `inf_bad` as `b=*`, because there is no number for how overfull.
%
% Three things had to be true together for this to work here.
%
% §873: the second pass *is* the final one when `\emergencystretch` is zero,
% which is INITEX's value and trip.tex's. We ran the second pass as if it could
% fail and then a third that scored by different rules — a flat cost §859 does
% not have, and a per-point charge for overflow that is ours alone.
%
% §851: an overfull line deactivates its node in *every* pass, the final one
% included. What the final pass may do is take the break anyway, not keep the
% node. Leaving it active kept the paragraph's start in the list to the very
% end, so the last break never looked like the sole one and never got its
% artificial demerits.
%
% §851 again: `if b>threshold then goto continue` holds in the final pass too.
% It is not more permissive about *looseness*; the only thing it relaxes is
% failure.
%
% The paragraph below cannot be set within `\tolerance` at this measure —
% `\pretolerance` is high enough that the first pass tries and fails, so all
% three passes' worth of behaviour is on show in one narration.
\catcode`\{=1 \catcode`\}=2
\tracingonline=1
\tracingparagraphs=1
\hsize=100pt \parindent=0pt
\pretolerance=100 \tolerance=200
\font\rip=cmr10 \rip

AAAAAAAAAAAAAAAAAAAAAAAAAAAA AAA\par

% §830, the first thing `try_break` does:
%
%     if abs(pi)>=inf_penalty then
%       if pi>0 then return {inhibited by infinite penalty}
%       else pi:=eject_penalty {this breakpoint will be forced}
%
% So every forcing penalty *is* `eject_penalty` from there on — in the demerits
% and in §856's `p=`. trip.tex writes `\penalty-1000000000` and Knuth's
% transcript reports `p=-10000`; we reported the number as written.
AAA\penalty-1000000000 AAA\par
\end
